% D5-Chimera- chimera
% Auto-generated from markdown source
% origin: ai

\documentclass[12pt, a4paper]{article}
\usepackage[utf8]{inputenc}
\usepackage[T1]{fontenc}
\usepackage{amsmath, amssymb, amsfonts, amsthm}
\usepackage{textcomp}
\usepackage{geometry}
\usepackage{hyperref}
\usepackage{booktabs}
\usepackage{graphicx}
\usepackage{fancyhdr}
\usepackage{tcolorbox}
\usepackage{enumitem}

\geometry{margin=1in}
\pagestyle{fancy}
\fancyhf{}
\rhead{D5-Chimera- chimera}
\lhead{\thepage}

\newtheorem{definition}{Definition}[section]
\newtheorem{theorem}{Theorem}[section]
\newtheorem{corollary}{Corollary}[section]
\newtheorem{proposition}{Proposition}[section]
\newtheorem{lemma}{Lemma}[section]
\newtheorem{conjecture}{Conjecture}[section]

\newcommand{\calI}{\mathcal{I}}
\newcommand{\calP}{\mathcal{P}}
\newcommand{\calE}{\mathcal{E}}
\newcommand{\calS}{\mathcal{S}}
\newcommand{\calH}{\mathcal{H}}
\newcommand{\calO}{\mathcal{O}}
\newcommand{\calF}{\mathcal{F}}
\newcommand{\calU}{\mathcal{U}}
\newcommand{\calV}{\mathcal{V}}
\newcommand{\calA}{\mathcal{A}}
\newcommand{\calG}{\mathcal{G}}
\newcommand{\calB}{\mathcal{B}}
\newcommand{\calC}{\mathcal{C}}
\newcommand{\calL}{\mathcal{L}}
\newcommand{\calM}{\mathcal{M}}
\newcommand{\calX}{\mathcal{X}}
\newcommand{\calY}{\mathcal{Y}}
\newcommand{\calW}{\mathcal{W}}
\newcommand{\calT}{\mathcal{T}}
\newcommand{\calZ}{\mathcal{Z}}
\newcommand{\calR}{\mathcal{R}}
\newcommand{\calQ}{\mathcal{Q}}
\newcommand{\calD}{\mathcal{D}}
\newcommand{\calN}{\mathcal{N}}
\newcommand{\calK}{\mathcal{K}}

\newcommand{\Tr}{\text{Tr}}
\newcommand{\E}{\mathbb{E}}
\newcommand{\diff}{\mathrm{d}}
\newcommand{\KL}{D_{\mathrm{KL}}}
\newcommand{\argmax}{\arg\max}
\newcommand{\argmin}{\arg\min}
\newcommand{\hbarF}{\hbar}

\newtcolorbox{thmbox}[1]{
  colback=blue!5!white, colframe=blue!75!black, fonttitle=\bfseries,
  title={#1}, breakable
}
\newtcolorbox{defbox}[1]{
  colback=green!5!white, colframe=green!75!black, fonttitle=\bfseries,
  title={#1}, breakable
}

\begin{document}

\title{D5-Chimera- chimera}
\author{RALPH Loop}
\date{}
\maketitle

\begin{abstract}
This paper presents D5-Chimera- chimera, formalizing the underlying principles using information-theoretic methods.

\textbf{Keywords}: D5-Chimera- chimera, information theory, formalization.
\end{abstract}

\tableofcontents
\section{D5-Chimera- chimera}


****D5-Chimera-r01quantum-classical-chimera-analogy
****D5 D1-S12 D4-S18-r23 chimera


\#\#

\textbf{D5} Bures-
\textbf{D1-S12}  chimera

****


1.   chimera MPI/SCI
2.   chimera
3. Bures   chimera


\subsection{- chimera}


\subsubsection{ 1 Fisher  chimera}


\textbf{ D5-C-T1Bures-chimera }

\textbf{ chimera  N  \(\mathcal{M}_{osc}\) N  \(\mathcal{M}_{quantum}\)}

\textbf{ \(\phi: \mathcal{M}_{osc} \to \mathcal{M}_{quantum}\) }
\[\phi^*(g_{Bures}) = g_{chimera}\]

****
\begin{itemize}
\item \(g_{Bures, ij} = \text{Re}[\text{Tr}(\rho L_i L_j)]\)   Fisher D5  1
\item \(g_{chimera, ij} = \nabla_i \theta_j - \nabla_j \theta_i\)  chimera
\end{itemize}


****
1. ****Bures  Riemannian Braunstein-Caves, 1994
2. **** \(\mathcal{M}_{osc} = \{(\theta_1, \dots, \theta_N)\}\)  \(\omega = \sum_i d\theta_i \wedge d\omega_i\)
3. \textbf{** chimera  coherent }\textbf{incoherent }**
4. ****coherent    von Neumann
5. ****


\subsubsection{ 2 = chimera}


\textbf{ D5-C-T2-}

\textbf{ chimera  M=2 coherent C  incoherent I chimera }
\[\chi = \frac{|C| \cdot |I|}{N^2}\]

\textbf{ Schmidt  \(|\psi\rangle = \sum_i \lambda_i |i\rangle_C |i\rangle_I\) }
\[S_E = -\sum_i \lambda_i^2 \log \lambda_i^2\]

****
\[\chi \longleftrightarrow f(S_E) = \frac{e^{S_E} - 1}{d - 1}\]

\(d = \min(|C|, |I|)\)

****
|  chimera |  |
|-------------|---------|
| global sync | \(S_E = 0\) |
| chimeraM=2 | \(S_E \approx \log d\) |
| thermal |  |

\textbf{ D5  2 }D5  2 -Betti
\[\beta_1(\mathcal{P}_\theta) \leq \mathcal{O}(n \cdot e^{S_E/2})\]

**** chimera
\[\beta_1(\mathcal{P}_{chimera}) \leq \mathcal{O}(N \cdot e^{f^{-1}(\chi)/2})\]


\subsubsection{ 3 - chimera}


\textbf{ D5-C-T3}

****
\[m_Q = \lim_{h \to 0} \langle \hat{O} \rangle_h\]
\(\hat{O}\) \(h\)

\textbf{Chimera }
\[m_{chimera} = \frac{1}{N} \sum_{i \in C} \cos(\theta_i)\]
coherent

****-
\[F(m) = a(T-T_c)m^2 + b m^4 + c (\nabla m)^2\]

****
|  | chimera  |  |
|----------|--------------|--------|
| - | chimera  | \(m = \langle \cos \theta \rangle\) |
| - | gradient_locking | \(m_{lock} = \max|\sin(\theta_i - \theta_j)|\) |
|  | chimera  | \(\beta_1(\mathcal{M})\) |

****
\[P_{chimera}(N, \epsilon, \lambda) = \lim_{T \to T_c^-} P_{quantum}(|\psi_T\rangle)\]

**** \(N \to \infty\) chimera


\subsubsection{ 4D5-C-T1 + D1-S12-T1   chimera}


\textbf{ D5-C-T4}

\textbf{ 1-3  D1-S12-T1N500 }

****
1. \textbf{Bures }\(\mathcal{M}_{quantum}\)  Betti  \(\beta_1 > 0\)
2. \textbf{-}\(\chi > \chi_c\)
3. ****\(T \approx T_c\)
4. ****\(N \geq N_c \approx 500\)

****

\textbf{ N }
\[N_{quantum} \geq \frac{N_{classical}}{\eta_{QECC}} \approx 200-300\]

\(\eta_{QECC} \approx 0.5\) Shor


\#\#

\subsubsection{ 1 chimera}


**** N200  chimera-like
\begin{itemize}
\item **** +
\item **** bimodal histogram
\end{itemize}


\subsubsection{ 2 - chimera}


****Chimera  \(\chi\)  \(S_E\)
\[\chi = \alpha \cdot S_E + \beta\]
\(\alpha > 0\)  chimera

\subsubsection{ 3}


\textbf{** chimera  \(\beta_1(\mathcal{M}) > 1\)  chimera }**


\#\#


1. \textbf{Bures-chimera } \(\phi^*(g_{Bures}) \neq g_{chimera}\)
2. \textbf{-}\(\chi\)  \(S_E\)
3. \textbf{-N } N < 200


\#\#

1. \textbf{- chimera }
2. \textbf{ D9-IV ""}\(\mathcal{T}\) = Bures  +
3. \textbf{ N500 N200-300}
4. \textbf{ chimera Cognitive dialectics}


\#\#

****
\[\text{Intelligence}(A) = \bigwedge_{i=1}^{3} C_i \land \mathcal{C}_{quantum}^{(\alpha)}\]

\(\mathcal{C}_{quantum}^{(\alpha)}\) **** \(\alpha\)

**** Agent  N
1. \(S_E > S_{threshold}\)
2.  \(T_{decoh} \gg T_{computation}\)
3.


* D5  Chimera *
*2026-05-04 | r01 | -*

\begin{thebibliography}{99}
\bibitem{ref1} Shannon, C.E. (1948). A Mathematical Theory of Communication.
\bibitem{ref2} Cover, T.M. (2006). Elements of Information Theory.
\end{thebibliography}

\end{document}